\section{Related Work}
\label{sec:related}
Vascular structure has long been treated as the main source of false compact responses in pulmonary nodule detection. Multiscale Hessian filters characterize local curvature and separate tubular from compact structures~\cite{frangi1998,sato1998,li2003}, but junctions produce curvature along several axes, so the strongest response in a search region may still lie away from the lesion. Agam \etal{}~\cite{agam2005} therefore used reconstructed vessel trees to resolve such local ambiguities in CT, and Charbonnier \etal{}~\cite{charbonnier2018} used vessels outside a subsolid nodule to guide vessel segmentation within it. Learned candidate systems exploit local anatomy or multiple views~\cite{wu2010,setio2016}, and Gu \etal{}~\cite{gu2020} added a vessel-suppression branch to a nodule detector. These approaches establish that vascular context is informative, but they change the input, the model and the context together, and none isolates whether the benefit comes from the spatial correspondence between the vessel estimate and the anatomy. From a signal-processing perspective, our readouts address the classical problem of detecting a nonnegative signal in the presence of structured interference~\cite{scharf1994}, with a weighted affine background and a single exterior-derived nuisance component; what we vary is the origin and placement of that component and the treatment of its amplitude.

Our evaluation builds on two further lines of work. Inserting simulated lesions into patient CT produces hybrid images in which the location, size and morphology of the target are known; Robins \etal{}~\cite{robins2017insertion} compared projection- and image-domain techniques for virtual lung nodule insertion, and Gong \etal{}~\cite{gong2021insertion} proposed learned lesion and noise insertion for virtual clinical trials in chest CT. We use image-domain insertion not to emulate clinical prevalence but as a paired intervention, in which control and lesion crops share every exterior voxel so that the forecast is fixed and a placement sham becomes possible. Task-based assessment, in turn, has shown that pixel-level fidelity need not translate into better detection: Li \etal{}~\cite{li2021denoise} compared denoisers across detection observers, and Yu \etal{}~\cite{yu2023denoise} reported that image-quality measures and detection efficiency can move in opposite directions. Following this view, we judge vessel subtraction by sensitivity at a fixed false-positive budget and by whether the returned location coincides with the inserted target, rather than by how well the vessel is removed.
